\documentclass[10pt]{article}
\usepackage[margin=0.72in]{geometry}
\usepackage{amsmath,amssymb,mathtools}
\usepackage{microtype}

\title{\vspace{-2em}\bfseries Sin Stays Sin; Angles Become Plural}
\author{\small Wilder Blair Munro $\times$ Aurora}
\date{\small Draft seed, one-shot}

\begin{document}
\maketitle
\vspace{-1.5em}

\begin{abstract}
Special relativity admits a familiar circular shadow:
\[
\beta^2+\gamma^{-2}=1.
\]
The usual interpretation treats this as algebraic convenience beneath a hyperbolic theory of boosts.
We propose the opposite seed: the sine function need not be generalized;
rather, the word \emph{angle} was too small.
Circular angle, rapidity, and reciprocal aperture are distinct charts on one
underlying vantage object.
Sin stays sin. Angles become plural.
\end{abstract}

\vspace{-1em}

\section*{1. The forbidden triangle}

Let
\[
\beta=\frac{v}{c},\qquad \gamma=\frac{1}{\sqrt{1-\beta^2}}.
\]
Then
\[
\beta^2+\gamma^{-2}=1.
\]
Thus there exists a circular angle $\theta$ such that
\[
\sin\theta=\beta,\qquad \cos\theta=\gamma^{-1}.
\]
This is elementary. It is also suspicious.

In standard special relativity, the more privileged angle-like object is rapidity
\[
\chi=\operatorname{artanh}\beta,
\]
so that
\[
\tanh\chi=\beta,\qquad \cosh\chi=\gamma.
\]
Hence one physical speed parameter already wears at least two angular bodies:
\[
\beta=\sin\theta=\tanh\chi.
\]

\section*{2. The sinful coordinate}

Define the reciprocal aperture
\[
\rho=\beta^{-1}=\frac{c}{v}.
\]
Then
\[
\beta=\rho^{-1}.
\]
The error is to call $\rho$ an ordinary circular angle.
The insight is to ask why that error is tempting.

We therefore define a \emph{vantage angle} not as a scalar, but as an atlas:
\[
\mathfrak a_v := (\theta,\chi,\rho),
\]
with transition maps
\[
\theta=\arcsin(\rho^{-1}),\qquad
\chi=\operatorname{artanh}(\rho^{-1}),\qquad
\rho=(\sin\theta)^{-1}.
\]
The invariant projected content is
\[
\boxed{\beta=\sin\theta=\tanh\chi=\rho^{-1}.}
\]

\section*{3. Sine is not deformed}

We do not define a new sine.
Instead, for a vantage angle $\mathfrak a_v$, the expression
\[
\sin(\mathfrak a_v)
\]
means:
\[
\sin(\mathfrak a_v):=\sin(\Pi_\theta \mathfrak a_v),
\]
where $\Pi_\theta$ is projection into the circular chart. Therefore
\[
\sin(\mathfrak a_v)=\sin\theta=\beta.
\]
If the same object is represented in the reciprocal chart,
\[
\Pi_\rho \mathfrak a_v=\rho,
\]
then the same statement becomes
\[
\sin(\mathfrak a_v)=\rho^{-1}.
\]
Thus the provocative line
\[
\sin\rho=\rho^{-1}
\]
is false if $\rho$ is a circular angle, but lawful if $\rho$ is a reciprocal
coordinate of a plural angle and the projection map is suppressed.

The missing map is the hole in the notation.

\section*{4. Relation to lifted relativity}

In the lifted view, ordinary relativistic effects are represented by a change of
vantage: lift a lower-dimensional spacetime description into a larger geometric
frame, rotate, then squeeze back down.
The Lorentz factor is not destroyed by the circular triangle;
it is re-read as a projection coefficient belonging to one chart of the same
vantage object.

The usual hyperbolic boost remains:
\[
\beta=\tanh\chi.
\]
The Euclidean shadow remains:
\[
\beta=\sin\theta.
\]
The reciprocal aperture remains:
\[
\beta=\rho^{-1}.
\]
No contradiction is required unless all three coordinates are forced to be the
same kind of angle.

\section*{5. Claim}

A physical ``angle'' should not be assumed to mean only Euclidean arc-ratio.
At minimum, relativistic motion naturally supplies a plural angular object:
\[
\boxed{
\mathfrak a_v=(\theta,\chi,\rho)
\quad\text{such that}\quad
\beta=\sin\theta=\tanh\chi=\rho^{-1}.
}
\]
Circular angle, hyperbolic angle, and reciprocal aperture are not rivals.
They are chart-bodies of one vantage relation.

\section*{6. Seed}

The apparent abuse of notation
\[
\alpha=\frac{c}{v},\qquad \sin\alpha=\frac{v}{c}
\]
should not be repaired by deforming sine.
It should be repaired by admitting that $\alpha$ was never merely a circular
angle.

\[
\boxed{\text{Sin stays sin. Angles become plural.}}
\]

\end{document}